\documentclass[12pt]{article}
\usepackage[a4paper,margin=1in]{geometry}
\usepackage[T1]{fontenc}
\usepackage{lmodern,microtype}
\usepackage{amsmath,amssymb,amsthm,mathtools}
\usepackage{booktabs,array,enumitem,xcolor}
\usepackage[numbers,sort&compress]{natbib}
\usepackage[colorlinks=true,linkcolor=blue!55!black,citecolor=blue!55!black,urlcolor=blue!55!black]{hyperref}
\linespread{1.07}

\newtheorem{theorem}{Theorem}[section]
\newtheorem{proposition}[theorem]{Proposition}
\newtheorem{corollary}[theorem]{Corollary}
\theoremstyle{remark}
\newtheorem{remark}[theorem]{Remark}
\newcommand{\norm}[1]{\left\lVert#1\right\rVert}

\title{Ten Layers from Orthogonal Clock Failure to an Oblique Feshbach Frontier\\
\large RH-182--RH-190 Route Review and the Next Physical Gate}
\author{Liang Wang}
\date{July 2026}

\begin{document}
\maketitle

\begin{abstract}
This paper reviews RH-182--RH-190, the next ten-layer exploration after the
reset-history/cycle review of RH-181.  The batch began with a strict test of
the data-derived orthogonal temporal clock at the predeclared lengths
$L=r-3$ and $L=r-4$.  The exact weighted-cycle construction was correct, but
all 126 windows failed the common wrap/primal/adjoint gate; even the best
projective endpoint return was too large.  This closes the simplest
orthogonal physical clock branch at the audited finite anchors.

The next layer used the physical observation to construct a balanced
biorthogonal source/observation clock.  The algebra is exact: the cross-Gram
SVD gives $W^*V=I$, the optimal oblique condition is the inverse smallest
cross singular value, and compressed spectra are gauge-covariant.  In the
physical audit, 12 of 38 $\sigma=0.01,L=4$ windows pass a two-sided $0.10$
residual gate; no length-three window passes.  This is a local floating
candidate, not a uniform realization.

The final layers quantify its boundary.  Cross-angle conditioning defeats a
coarse maximum-residual gate and every singular-value clipping regularizer;
scalar gauge balancing leaves the sharp directional product invariant, with
eight windows having absolute coupling product below one.  An exact oblique
Feshbach determinant factorization is proved and audited.  The elementary
norm-only complement-resolvent budget then fails on all 126 windows, even
under an optimistic unit complement factor.

The updated frontier is therefore precise:
\begin{equation*}
 \begin{gathered}
 \text{local biorthogonal packet candidate}\\
 \downarrow\\
 \text{validated physical complement contours }D\\
 \downarrow\\
 \text{uniform margin }K,\ \text{transport }H,\ \text{cloud ledger }Q
 \end{gathered}
\end{equation*}
Gate A and all Hilbert--Polya downstream gates remain open.  No zero
identification or Riemann-hypothesis conclusion is claimed.
\end{abstract}

\section{Coordinate before this batch}

RH-181 left two architecture-relative routes.  The reset-history branch had
an exact memory-to-history realization but no history-to-transfer map.  The
finite-cycle branch had exact determinant algebra but no physical calibration
or cycle-to-transfer map \citep{WangRH181}.  The recommended next experiment
was a target-independent finite temporal clock using $L=r-3$ and $L=r-4$.

The present review records not only what survived, but also which negative
shortcuts are now closed.  This matters because a local positive numerical
window is not the same object as an all-level physical leaf.

\section{Layer-by-layer ledger}

\begin{center}
\small
\begin{tabular}{p{0.09\textwidth}p{0.61\textwidth}p{0.18\textwidth}}
\toprule
paper & principal result & strict status\\
\midrule
RH-182 & finite normalized orbit clock, exact rank-one wrap, 126-window audit
& orthogonal branch rejected\\
RH-183 & optimal phase/scalar wrap and projective lower bound
& finite obstruction proved\\
RH-184 & balanced cross-Gram biorthogonal frames and optimal conditioning
& finite algebra proved\\
RH-185 & physical source/observation bi-Krylov calibration
& local $L=4$ candidate\\
RH-186 & oblique conditioning amplification
& coarse gate rejected\\
RH-187 & exact clipping Pareto tradeoff
& clipping branch rejected\\
RH-188 & scalar gauge balance and invariant directed product
& partial positive signal\\
RH-189 & exact oblique Feshbach determinant identity
& finite algebra proved\\
RH-190 & norm-only complement resolvent budget
& norm route rejected\\
\bottomrule
\end{tabular}
\end{center}

The machine ledger contains 2,960 finite cases, windows, or sweep items and
zero recorded formula/identity failures.  This is an implementation ledger,
not a statistical sample and not evidence that the open physical leaves are
true.

\section{The orthogonal branch is genuinely closed at finite level}

For a normalized source orbit $x_j$ and temporal synthesis $J$, RH-182 pairs
the chain with a weighted cycle.  The intertwining error is exactly
\begin{equation}\label{eq:rank-one}
 \mathcal AJ-JC
 =a_{t+L-1}(x_{t+L}-\omega x_t)e_{L-1}^*.
\end{equation}
The cycle spectrum is an exact root grid.  Thus the test did not fail because
of a poorly assembled determinant.

RH-183 proves that even optimizing over all phases and all final-edge
scalars leaves the projective lower bound
\begin{equation}\label{eq:projective}
 a_{t+L-1}\sqrt{1-|\langle x_t,x_{t+L}\rangle|^2}.
\end{equation}
Across 126 windows, no projective return is at most $0.25$.  The minimum
adjoint residual in the polar orthogonal construction is $0.7143$.  Hence
the following finite branch is closed:
\begin{equation}\label{eq:closed-orthogonal}
 \text{one orthogonal temporal packet}
 +\text{one phase/scalar wrap}
 \Longrightarrow\text{two-sided physical clock}.
\end{equation}

This does not reject a different temporal span, a lagged map, or two distinct
right/left spaces.

\section{The biorthogonal algebraic bridge}

RH-184 starts with right and left orthonormal temporal bases $Q_R,Q_L$ and
cross Gram $H=Q_L^*Q_R$.  If
$H=U\Sigma V^*$, the balanced frames are
\begin{equation}\label{eq:balanced}
 V_R=Q_RV\Sigma^{-1/2},
 \qquad W_L=Q_LU\Sigma^{-1/2},
 \qquad W_L^*V_R=I.
\end{equation}
The exact optimal conditioning is
\begin{equation}\label{eq:optimal-conditioning}
 \norm{V_RW_L^*}=\sigma_{\min}(H)^{-1}.
\end{equation}
This is not a frame-choice artifact.  Every biorthogonal realization on the
same two subspaces has norm product at least this value.

The compressed block and directed residuals are
\begin{equation}\label{eq:directed}
 K=W_L^*AV_R,
 \qquad
 R_R=AV_R-V_RK,
 \qquad
 R_L=A^*W_L-W_LK^*.
\end{equation}
This is the exact typed replacement for an orthogonal reducing packet.

\section{Local physical calibration}

RH-185 constructs the right history from $A^jS$ and the left history from
$(A^*)^jO^*$.  The same 126 windows are tested.  The results are:
\begin{center}
\begin{tabular}{lrrr}
\toprule
branch & windows & two-sided passes & minimum condition\\
\midrule
all length-three candidates & 88 & 0 & $48.22$\\
$\sigma=0.01,L=4$, left & 19 & 5 & $88.11$\\
$\sigma=0.01,L=4$, right & 19 & 7 & $179.75$\\
\bottomrule
\end{tabular}
\end{center}
The accepted length-four windows have phase-grid errors near $0.1$ radians
and radial errors near a few hundredths.  The best residual pair is
$(0.02332,0.02498)$.

This selects $L=4$ as a local calibration at the only actual RH-15/RH-151
overlap.  It does not select it at untested scales and does not prove a cloud
degree law.

\section{Conditioning and regularization}

RH-186 inserts $\chi=\sigma_{\min}(H)^{-1}$ into a coarse coordinate gate
\begin{equation}\label{eq:amp}
 \chi\max(\epsilon_R,\epsilon_L)<1.
\end{equation}
The minimum amplified residual over all 126 windows is $10.253$, so the gate
has zero successes.  RH-187 then clips the cross-Gram singular values.  If
$\gamma=\sigma_{\min}(H)$, residual level $\epsilon$, and clipping level
$\tau$, the combined budget is exactly
\begin{equation}\label{eq:clip}
 1+\frac{\epsilon-\gamma}{\max(\gamma,\tau)}
 \quad (\tau\ge\gamma).
\end{equation}
Strict contraction exists if and only if $\epsilon<\gamma$.  The physical
minimum ratio is again $10.253$; 1,638 clipping sweeps yield zero strict
successes.  This closes the simple spectral-clipping repair.

\section{The directional product remains alive}

RH-188 separates the two couplings.  Under
$V\mapsto\alpha V$, $W\mapsto\alpha^{-1}W$, the couplings transform as
$c\mapsto\alpha c$ and $b\mapsto\alpha^{-1}b$.  Hence $bc$ is invariant and
the maximum is minimized at $\sqrt{bc}$.

The physical absolute product has minimum $0.2486$ and eight windows have
$bc<1$, all in the local length-four branch.  Twelve local windows have
relative product below $0.01$.  Therefore the sharp Schur route is not
eliminated by the maximum-norm negative result.  It still needs packet and
complement resolvent factors.

\section{Exact Feshbach structure}

RH-189 makes the missing factors type-correct.  Extend $V$ by a basis $Z$ of
$\ker W^*$ and write $S=[V,Z]$.  Then
\begin{equation}\label{eq:block}
 S^{-1}AS=\begin{pmatrix}K&B\\C&D\end{pmatrix},
\end{equation}
and
\begin{equation}\label{eq:feshbach}
 \det(zI-A)=\det(zI-D)
 \det\left(zI-K-B(zI-D)^{-1}C\right).
\end{equation}
The 240-case identity audit has maximum relative determinant error
$1.85\times10^{-11}$.  The physical complement resolvent is now a specific
operator, not an informal remainder.

\section{The norm-only complement route fails}

RH-190 applies the universal estimate
\begin{equation}\label{eq:norm-only}
 \norm{D}\le\chi\norm A
\end{equation}
in the orthonormal complement coordinates of RH-189.
With contour radius $0.4$ of the cycle half-spacing, every one of 126
windows has negative clearance.  The smallest complement norm bound is
$71.46$, while the largest contour minimum modulus is only
$0.7188$.  Even assigning the complement resolvent the artificial value one
leaves zero Schur successes; the optimistic product has minimum $1.1116$.

This is a coarse norm failure, not a spectral no-go.  It tells us exactly
where the next computation belongs: validated sample inverses on the actual
contours, with mesh covering and operator balls as in RH-167--168.

\section{Updated frontier}

The current finite route can be written
\begin{equation}\label{eq:route}
 \begin{aligned}
 &\text{balanced bi-Krylov algebra [proved]}\\
 &\quad+\ \text{local }(\sigma=0.01,L=4)\text{ candidate [floating]}\\
 &\quad\longrightarrow\ D_{\rm phys}\text{ [validated contours, open]}\\
 &\quad\longrightarrow\ K_{\rm phys},H_{\rm phys},Q\text{ [open]}.
 \end{aligned}
\end{equation}
The reset-history branch remains open independently.  The orthogonal clock
shortcut is removed, but the biorthogonal branch is not a proof of physical
interface R.

\section{Decision tree at the present frontier}

The next contour calculation has two informative outcomes.  If every
predeclared local window fails because the complement contour intersects the
validated pseudospectrum or because the full Schur product is at least one,
then the present $L=4$ realization is rejected at those finite anchors.  The
reset-history branch and other biorthogonal seeds remain logically open, but
this packet choice should not be carried forward.

If at least one window has a positive continuous margin, the conclusion is
still finite: RH-189 then identifies the full Riesz count as
\begin{equation}\label{eq:count-ledger}
 N_A(\Gamma)=N_K(\Gamma)+N_D(\Gamma).
\end{equation}
The complement count must be shown to vanish, or retained explicitly.  Only
after this count is stable under outward operator balls can the route move to
uniform margin $K$ and shell transport $H$.

This decision tree makes the present exploration falsifiable.  A local
candidate is not protected by the earlier phase agreement; it survives only
if the exact complement calculation does.

\section{Position inside the five macro gates}

The entire batch remains inside Gate A of the larger roadmap.  More
precisely, it concerns the physical realization and finite Riesz interface
that precede an intrinsic cloud ledger.  None of the following downstream
tasks has begun here:
\begin{enumerate}[leftmargin=2em]
\item completion of the non-self-adjoint dynamics into a time-oriented
 unitary or scattering object;
\item construction of a self-adjoint generator;
\item derivation of a $T\log T$ counting law from that generator;
\item derivation of von Mangoldt prime-power weights from a trace formula;
\item equality of the resulting spectral divisor with the completed zeta
 divisor.
\end{enumerate}
Keeping this macro position explicit prevents finite packet eigenvalues from
being described as zeta zeros or Hilbert--Polya eigenvalues.

\section{Batch reproducibility and evidence classes}

The 2,960-item aggregate combines formula trials, physical windows, and
regularization sweep points.  These objects have different logical types.
Formula trials test exact finite identities against independent matrix
computations.  Physical windows evaluate declared deterministic gates.
Threshold sweeps map a branch-specific Pareto curve.  Their counts are added
only for archive completeness; they are not exchangeable observations and
do not define a statistical confidence level.

Each paper therefore carries three separate records: a theorem ledger, a
machine result, and a roadmap boundary.  The batch manifest hashes the
publication files across RH-182--RH-191.  This makes later revisions able to
distinguish a changed theorem, a changed dataset, and a changed
interpretation rather than treating the ten layers as one opaque numerical
experiment.

\section{Final boundary}

RH-182--190 prove finite weighted-clock, projective, biorthogonal, gauge,
Feshbach, and norm-budget statements and provide reproducible finite audits.
They do not prove uniform cross-angle control, a physical complement inverse,
continuous Schur margins, Riesz shell ranks, shell transport, cloud ledger Q,
complement limit U, canonicity Z, marked limit T, or macro Gate A.

Gates B--E, a self-adjoint generator, the $T\log T$ law, a von Mangoldt
trace formula, a zeta divisor identity, Hilbert--Polya, and the Riemann
Hypothesis remain untouched.

\bibliographystyle{plainnat}
\bibliography{references}
\end{document}
